\section{Related Work}
\label{sec:related}

\para{Neural thickets and weight-space ensembles.}
\citet{gan2026neural} discover that pretrained LLM weights are surrounded by dense clusters of task-specific expert solutions, reachable via random Gaussian perturbations. Their \randopt algorithm selects top-$K$ perturbed models by validation accuracy and ensembles predictions via majority vote, achieving performance competitive with PPO and GRPO on reasoning benchmarks at $O(1)$ training cost. The same research group previously explored task-aware ensemble construction through \textsc{Split-Ensemble} \citep{chen2024split}, which splits classification into complementary subtasks for OOD-aware training. Our work extends neural thickets to the vision domain and provides the first evaluation of their adversarial robustness and OOD detection capabilities.

\para{Model soups and weight interpolation.}
A related line of work explores linear combinations of model weights for robustness. \citet{croce2023seasoning} show that ``model soups'' --- convex or extrapolated combinations of adversarially fine-tuned models --- can smoothly trade off robustness to different $\ell_p$-norm adversaries without retraining. \citet{rame2022diverse} demonstrate that averaging weights of diverse models improves OOD generalization. These approaches operate in the same weight-space paradigm as neural thickets, but use structured directions (fine-tuned models) rather than random perturbations. Our work reveals why the random perturbation approach fails for vision: the perturbation must be orders of magnitude smaller than the structured directions used in model soups.

\para{Ensemble diversity for adversarial robustness.}
The connection between ensemble diversity and adversarial robustness has been extensively studied. \citet{pang2019improving} propose Adaptive Diversity Promoting (ADP) regularization, which encourages non-maximal predictions to be mutually orthogonal across ensemble members, reducing adversarial transferability. \citet{zhao2024ensemble} integrate low-curvature models for ensemble defense, showing that diversity combined with flat loss landscapes improves robustness. \citet{hu2024meat} propose median-based weight averaging for self-ensemble adversarial training. A common finding across these works is that diversity must be explicitly encouraged during training --- incidental diversity from different random seeds or perturbations is typically insufficient \citep{fort2019deep}. Our experiments confirm this finding in the extreme case of weight-perturbation ensembles.

\para{OOD detection with ensembles.}
Deep ensembles are a natural tool for OOD detection, as ensemble disagreement can signal distributional shift. However, \citet{xia2022usefulness} demonstrate that ensemble diversity (measured by mutual information) is not reliably useful for OOD detection at scale; instead, averaging task-specific scores (\eg energy) over ensemble members is more effective. Our OOD detection results corroborate this finding: energy-based detection consistently outperforms mutual information across all thicket variants.

\para{Loss landscape geometry and robustness.}
\citet{fort2019deep} show that deep ensembles derive their uncertainty estimates from exploring different modes of the loss landscape, and that the quality of uncertainty quantification depends on mode diversity rather than proximity in weight space. \citet{lu2024sharpbalance} identify a tradeoff between sharpness of individual learners and ensemble diversity. These geometric perspectives help explain our negative result: perturbations within a single sharp basin of the loss landscape cannot access the diverse modes needed for robustness.
